\chapter{Open-Source Softwares}
\label{chap:software-data}

% ============================================================
% SOFTWARE OVERVIEW
% ============================================================

\section{Software Overview}

\begin{table}[htbp]
\centering
\caption{Comparison of open-source software for EAGLE requirements}
\label{tab:eagle_software}
\renewcommand{\arraystretch}{1.25}
\setlength{\tabcolsep}{3pt}

\begin{adjustbox}{max width=\textwidth}
\begin{tabular}{
>{\raggedright\arraybackslash}p{4.2cm}
>{\raggedright\arraybackslash}p{2.6cm}
>{\raggedright\arraybackslash}p{2.6cm}
>{\raggedright\arraybackslash}p{2.6cm}
>{\raggedright\arraybackslash}p{2.6cm}
>{\raggedright\arraybackslash}p{2.6cm}}
\hline

\textbf{EAGLE Requirement} &
\textbf{PyPSA} &
\textbf{pandapower} &
\textbf{OpenDSS} &
\textbf{MATPOWER} &
\textbf{VeraGrid} \\

\hline

% ============================================================
% GENERAL
% ============================================================

\textbf{Open-source} &
\cellcolor{green!25} Yes &
\cellcolor{green!25} Yes &
\cellcolor{green!25} Yes &
\cellcolor{green!25} Yes &
\cellcolor{green!25} Yes \\

\textbf{Python-native workflow} &
\cellcolor{green!25} Native &
\cellcolor{green!25} Native &
\cellcolor{orange!25} Via interfaces &
\cellcolor{red!25} MATLAB / Octave &
\cellcolor{green!25} Native \\

\textbf{Programmatic automation} &
\cellcolor{green!25} Yes &
\cellcolor{green!25} Yes &
\cellcolor{green!25} Yes &
\cellcolor{green!25} Yes &
\cellcolor{green!25} Yes \\


% ============================================================
% NETWORK MODELLING
% ============================================================

\textbf{HV transmission modelling} &
\cellcolor{green!25} Strong &
\cellcolor{green!25} Strong &
\cellcolor{orange!25} Secondary focus &
\cellcolor{green!25} Strong &
\cellcolor{green!25} Strong \\

\textbf{MV / distribution modelling} &
\cellcolor{orange!25} Possible &
\cellcolor{green!25} Strong &
\cellcolor{green!25} Strong &
\cellcolor{orange!25} Possible &
\cellcolor{green!25} Strong \\

\textbf{HVDC modelling} &
\cellcolor{green!25} Yes &
\cellcolor{green!25} Yes &
\cellcolor{orange!25} Limited / indirect &
\cellcolor{orange!25} Limited &
\cellcolor{green!25} Yes \\


% ============================================================
% ELECTRICAL ANALYSIS
% ============================================================

\textbf{Non-linear AC power flow} &
\cellcolor{green!25} Yes &
\cellcolor{green!25} Yes &
\cellcolor{green!25} Yes &
\cellcolor{green!25} Yes &
\cellcolor{green!25} Yes \\

\textbf{DC / linear power flow} &
\cellcolor{green!25} Yes &
\cellcolor{green!25} Yes &
\cellcolor{orange!25} Limited &
\cellcolor{green!25} Yes &
\cellcolor{green!25} Yes \\

\textbf{Line loading / thermal limits} &
\cellcolor{green!25} Yes &
\cellcolor{green!25} Yes &
\cellcolor{green!25} Yes &
\cellcolor{green!25} Yes &
\cellcolor{green!25} Yes \\

\textbf{Voltage constraints} &
\cellcolor{green!25} Yes &
\cellcolor{green!25} Yes &
\cellcolor{green!25} Yes &
\cellcolor{green!25} Yes &
\cellcolor{green!25} Yes \\

\textbf{Time-series analysis} &
\cellcolor{green!25} Strong &
\cellcolor{green!25} Strong &
\cellcolor{green!25} Strong &
\cellcolor{orange!25} MOST / additional workflow &
\cellcolor{green!25} Yes \\

\textbf{Contingency / N-1 analysis} &
\cellcolor{green!25} Yes &
\cellcolor{green!25} Yes &
\cellcolor{orange!25} Possible &
\cellcolor{orange!25} Additional workflow &
\cellcolor{green!25} Yes \\

\textbf{PTDF / LODF analysis} &
\cellcolor{orange!25} Possible &
\cellcolor{orange!25} Possible &
\cellcolor{red!25} Not core &
\cellcolor{orange!25} Possible &
\cellcolor{green!25} Native \\


% ============================================================
% OPTIMISATION / CAPACITY
% ============================================================

\textbf{Optimal Power Flow (OPF)} &
\cellcolor{green!25} Strong &
\cellcolor{green!25} Yes &
\cellcolor{orange!25} Limited &
\cellcolor{green!25} Strong &
\cellcolor{green!25} Yes \\

\textbf{Long-term planning / expansion} &
\cellcolor{green!25} Strong &
\cellcolor{orange!25} Limited &
\cellcolor{orange!25} Limited &
\cellcolor{orange!25} Additional tools &
\cellcolor{orange!25} Possible \\

\textbf{Hosting-capacity studies} &
\cellcolor{orange!25} Custom methodology &
\cellcolor{orange!25} Custom methodology &
\cellcolor{green!25} Strong for distribution &
\cellcolor{orange!25} Custom methodology &
\cellcolor{green!25} Nodal functionality \\

\textbf{Large-scale network studies} &
\cellcolor{green!25} Strong &
\cellcolor{green!25} Yes &
\cellcolor{orange!25} Distribution-oriented &
\cellcolor{green!25} Strong &
\cellcolor{green!25} Yes \\


% ============================================================
% DATA / INTEROPERABILITY
% ============================================================

\textbf{PyPSA-Eur compatibility} &
\cellcolor{green!25} Native &
\cellcolor{orange!25} Conversion required &
\cellcolor{orange!25} Conversion required &
\cellcolor{orange!25} Conversion required &
\cellcolor{green!25} PyPSA import \\

\textbf{CIM / CGMES interoperability} &
\cellcolor{orange!25} External workflow &
\cellcolor{orange!25} Partial / conversion &
\cellcolor{orange!25} External workflow &
\cellcolor{orange!25} External workflow &
\cellcolor{green!25} Supported \\


% ============================================================
% GIS / EAGLE INTEGRATION
% ============================================================

\textbf{Geographical network data} &
\cellcolor{green!25} Coordinates &
\cellcolor{green!25} Geodata &
\cellcolor{orange!25} Basic coordinates &
\cellcolor{red!25} Not GIS-oriented &
\cellcolor{green!25} Geographic data \\

\textbf{GIS / map visualisation} &
\cellcolor{orange!25} Basic / external &
\cellcolor{green!25} Yes &
\cellcolor{orange!25} External &
\cellcolor{red!25} No &
\cellcolor{green!25} Yes \\

\textbf{Python GIS / SiteOptimizer integration} &
\cellcolor{green!25} Direct &
\cellcolor{green!25} Direct &
\cellcolor{orange!25} Via Python interface &
\cellcolor{red!25} Additional integration &
\cellcolor{green!25} Direct \\

\hline

\textbf{Main strength for EAGLE} &
\cellcolor{green!25}\textbf{Large-scale planning and optimisation} &
\cellcolor{green!25}\textbf{Detailed electrical analysis} &
\cellcolor{green!25}\textbf{Distribution and hosting capacity} &
\cellcolor{green!25}\textbf{Power flow and OPF reference} &
\cellcolor{green!25}\textbf{Transmission analysis and interoperability} \\

\hline

\end{tabular}
\end{adjustbox}

\vspace{0.2cm}

\begin{minipage}{0.95\textwidth}
\footnotesize

Green indicates that the capability is directly supported and well aligned
with the corresponding EAGLE requirement. Orange indicates that the
capability is available with limitations, through additional modules or
interfaces, or is not a primary focus of the software. Red indicates that
the capability is not natively supported or is poorly aligned with the
requirement.

Hosting capacity is not necessarily a single software function. In several
frameworks it must be implemented as a methodology combining power-flow or
optimal-power-flow calculations with voltage, thermal and operational
constraints.

\end{minipage}

\end{table}



\section{Software Description}

% ============================================================
% PyPSA
% ============================================================

\subsection{PyPSA}

\begin{description}

    \item[\textbf{Developer / organisation:}]
    PyPSA community. Development and maintenance are currently led by the
    Department of Digital Transformation in Energy Systems at the Technical
    University of Berlin.

    \item[\textbf{Access:}] \hfill
    \begin{itemize}
        \item Official website:
        \url{https://pypsa.org/}

        \item Documentation:
        \url{https://docs.pypsa.org/}

        \item Source repository:
        \url{https://github.com/PyPSA/PyPSA}
    \end{itemize}

    \item[\textbf{Description:}]
    PyPSA (Python for Power System Analysis) is an open-source Python
    framework for modelling, simulating and optimising electrical and
    multi-energy systems. It is particularly oriented towards large-scale
    transmission systems, renewable integration, storage, sector coupling
    and operational and investment optimisation over long time series.

    \item[\textbf{Main capabilities:}] \hfill
    \begin{itemize}
        \item Static non-linear AC power flow.
        \item Linearised power flow.
        \item Economic dispatch.
        \item Linear optimal power flow.
        \item Security-constrained optimal power flow.
        \item Unit commitment.
        \item Generation and storage capacity optimisation.
        \item Transmission expansion optimisation.
        \item Multi-period and multi-investment-period optimisation.
        \item Renewable generation modelling.
        \item Storage and hydro modelling.
        \item Sector coupling and multi-carrier energy systems.
        \item Network and result visualisation.
    \end{itemize}

    \item[\textbf{Network components:}] \hfill
    \begin{itemize}
        \item Buses.
        \item AC transmission lines.
        \item Transformers.
        \item Controllable links and HVDC connections.
        \item Generators.
        \item Loads.
        \item Storage units.
        \item Energy stores.
        \item Shunt impedances.
        \item Multi-carrier conversion technologies.
    \end{itemize}

    \item[\textbf{Power-system analysis:}] \hfill
    \begin{itemize}
        \item Full non-linear AC power flow using Newton--Raphson.
        \item Linearised AC/DC network analysis.
        \item Network-constrained economic dispatch.
        \item Linear optimal power flow.
        \item Security-constrained linear optimal power flow considering
        network contingencies.
        \item Operational optimisation over multiple time steps.
        \item Capacity-expansion and transmission-expansion optimisation.
    \end{itemize}

    \item[\textbf{Data handling and automation:}]
    PyPSA is implemented in Python and stores network components and
    time-dependent information in structures closely integrated with pandas.
    This facilitates automated workflows, modification of network parameters,
    batch simulations, processing of long time series and integration with
    external Python data-analysis libraries.

    \item[\textbf{GIS / spatial capabilities:}]
    Network buses can contain geographical coordinates and PyPSA provides
    network plotting and exploration functionality. Its Python-based
    architecture facilitates integration with GeoPandas and other GIS
    libraries. The PyPSA ecosystem also includes PyPSA-Eur, which provides
    geographically referenced European energy-system models.

    \item[\textbf{Input / output formats:}]
    PyPSA networks can be imported and exported using machine-readable
    formats including NetCDF, CSV and HDF5-based workflows depending on the
    functionality used. Network and result tables can also be directly
    accessed and processed through Python and pandas.

    \item[\textbf{Main advantages for EAGLE:}] \hfill
    \begin{itemize}
        \item Native Python environment suitable for automated EAGLE
        workflows.
        \item Direct compatibility with PyPSA-Eur, currently considered as
        a candidate source for the initial Spanish transmission-network model.
        \item Designed for large transmission networks and long time series.
        \item Strong support for renewable generation and storage.
        \item Network constraints can be included directly in optimisation.
        \item Suitable for studying future electricity-system scenarios.
        \item Supports both operational and infrastructure-investment
        optimisation.
        \item Results can be processed programmatically for integration with
        the spatial SiteOptimizer workflow.
    \end{itemize}

    \item[\textbf{Main limitations for EAGLE:}] \hfill
    \begin{itemize}
        \item Many optimisation functions rely on linearised network
        formulations rather than full non-linear AC equations.
        \item It is primarily oriented towards system-level and planning
        studies rather than highly detailed equipment-level electrical
        analysis.
        \item Detailed short-circuit and protection studies are not its
        primary purpose.
        \item Additional validation may be required when using reconstructed
        or estimated network parameters from PyPSA-Eur.
    \end{itemize}

    \item[\textbf{Potential application in EAGLE:}]
    PyPSA could serve as the main simulation and optimisation framework of
    the EAGLE electrical-network model. A Spanish network extracted from
    PyPSA-Eur could be enriched with official Spanish data and used to
    calculate network flows, line loading, congestion and available network
    capacity under different demand and generation scenarios. Its
    optimisation capabilities are particularly relevant for evaluating
    future data-centre demand, renewable integration and possible network
    reinforcement scenarios.

\end{description}


% ============================================================
% pandapower
% ============================================================

\subsection{pandapower}

\begin{description}

    \item[\textbf{Developer / organisation:}]
    Joint development involving the Energy Management and Power System
    Operation research group at the University of Kassel and the Grid
    Planning and Grid Operation Division at Fraunhofer IEE.

    \item[\textbf{Access:}] \hfill
    \begin{itemize}
        \item Official website:
        \url{https://www.pandapower.org/}

        \item Documentation:
        \url{https://pandapower.readthedocs.io/}

        \item Source repository:
        \url{https://github.com/e2nIEE/pandapower}
    \end{itemize}

    \item[\textbf{Description:}]
    pandapower is an open-source Python framework for modelling, analysis
    and optimisation of electrical power systems. It combines detailed
    element-based electrical models with a pandas-based data structure and
    is designed for highly automated static and quasi-static power-system
    studies.

    \item[\textbf{Main capabilities:}] \hfill
    \begin{itemize}
        \item AC power flow.
        \item DC power flow.
        \item Optimal power flow.
        \item State estimation.
        \item Short-circuit calculation.
        \item Topological network analysis.
        \item Time-series simulations.
        \item Controller-based simulations.
        \item Network plotting and geographical visualisation.
        \item Conversion from and to several external network formats.
    \end{itemize}

    \item[\textbf{Network components:}] \hfill
    \begin{itemize}
        \item Buses.
        \item AC lines.
        \item DC lines.
        \item Two-winding and three-winding transformers.
        \item Loads.
        \item Generators.
        \item Static generators.
        \item External grids.
        \item Switches.
        \item Shunts.
        \item Storage systems.
        \item Impedances and other electrical elements.
    \end{itemize}

    \item[\textbf{Power-system analysis:}] \hfill
    \begin{itemize}
        \item Newton--Raphson AC power flow.
        \item DC power flow.
        \item Optimal power flow.
        \item State estimation and bad-data analysis.
        \item IEC 60909 short-circuit calculations.
        \item Network topology analysis using graph-based methods.
        \item Quasi-static time-series simulation.
    \end{itemize}

    \item[\textbf{Data handling and automation:}]
    Network elements and simulation results are stored in pandas DataFrames,
    such as buses, lines, transformers, loads and their corresponding result
    tables. This structure is particularly convenient for automated data
    processing, filtering, modification, validation and integration with
    other Python libraries.

    \item[\textbf{GIS / spatial capabilities:}]
    pandapower supports geographical information associated with network
    elements and provides plotting functionality. Its Python environment
    also facilitates integration with external GIS libraries for more
    advanced spatial analysis.

    \item[\textbf{Input / output formats:}]
    Networks can be stored and exchanged through several formats and
    converters, including JSON, Excel and MATPOWER/PYPOWER-compatible
    representations. Network tables can also be directly manipulated and
    exported using pandas.

    \item[\textbf{Main advantages for EAGLE:}] \hfill
    \begin{itemize}
        \item Native Python integration.
        \item Detailed and validated electrical component models.
        \item Convenient pandas-based representation of network elements
        and results.
        \item Detailed calculation of bus voltages, branch flows, loading
        and electrical losses.
        \item State-estimation and short-circuit functionality if required
        in later EAGLE developments.
        \item NetworkX-based topology analysis can be useful for validating
        reconstructed network connectivity.
        \item Standard line and transformer libraries can help complete
        missing electrical parameters.
        \item Suitable for automated batch and time-series studies.
    \end{itemize}

    \item[\textbf{Main limitations for EAGLE:}] \hfill
    \begin{itemize}
        \item Long-term generation and transmission capacity-expansion
        optimisation is not its main design focus.
        \item Energy-system planning and sector coupling are less central
        than in PyPSA.
        \item Direct integration with PyPSA-Eur would require conversion of
        the network representation.
        \item For very large multi-period optimisation problems, a framework
        specifically designed for energy-system optimisation may be more
        convenient.
    \end{itemize}

    \item[\textbf{Potential application in EAGLE:}]
    pandapower could be used as a detailed electrical-analysis and validation
    framework for the EAGLE network. It is particularly suitable for checking
    AC power-flow results, voltage limits, line loading, electrical losses
    and network connectivity. It could therefore either serve as the main
    electrical simulation engine or complement a PyPSA-based planning model
    with more detailed electrical-network analyses.

\end{description}


% ============================================================
% OpenDSS
% ============================================================

\subsection{OpenDSS}

\begin{description}

    \item[\textbf{Developer / organisation:}]
    Electric Power Research Institute (EPRI).

    \item[\textbf{Access:}] \hfill
    \begin{itemize}
        \item Official documentation:
        \url{https://opendss.epri.com/}

        \item Official downloads:
        \url{https://sourceforge.net/projects/electricdss/}

        \item Python interface ecosystem:
        \url{https://dss-extensions.org/}

        \item OpenDSSDirect.py:
        \url{https://github.com/dss-extensions/OpenDSSDirect.py}
    \end{itemize}

    \item[\textbf{Description:}]
    OpenDSS (Open Distribution System Simulator) is an open-source electrical
    system simulation tool developed primarily for utility distribution
    systems and distributed-energy-resource integration. It supports detailed
    multi-phase network modelling and is particularly oriented towards
    distribution planning, DER studies and quasi-static time-series
    simulation.

    \item[\textbf{Main capabilities:}] \hfill
    \begin{itemize}
        \item Multi-phase power-flow analysis.
        \item Unbalanced distribution-network modelling.
        \item Quasi-static time-series simulation.
        \item Daily and yearly load-shape simulation.
        \item Distributed-generation and DER analysis.
        \item Voltage-regulation studies.
        \item Harmonic analysis.
        \item Fault and short-circuit studies.
        \item Energy-metering and loss analysis.
        \item Distribution planning studies.
    \end{itemize}

    \item[\textbf{Network components:}] \hfill
    \begin{itemize}
        \item Buses and nodes.
        \item Multi-phase lines and cables.
        \item Transformers.
        \item Loads.
        \item Generators.
        \item PV systems.
        \item Storage systems.
        \item Capacitors.
        \item Reactors.
        \item Voltage regulators.
        \item Protection and switching devices.
    \end{itemize}

    \item[\textbf{Power-system analysis:}] \hfill
    \begin{itemize}
        \item Multi-phase AC circuit analysis.
        \item Unbalanced power flow.
        \item Long quasi-static time-series simulations.
        \item Voltage and overload analysis.
        \item Harmonic studies.
        \item Fault and short-circuit calculations.
        \item Detailed analysis of distributed generation and storage.
    \end{itemize}

    \item[\textbf{Data handling and automation:}]
    OpenDSS uses its own text-based scripting language and can be automated
    through external programming interfaces. Python integration is possible
    through projects such as DSS-Python and OpenDSSDirect.py. This enables
    automated simulations and integration with the wider Python ecosystem,
    although these interfaces should be distinguished from the official
    EPRI implementation.

    \item[\textbf{GIS / spatial capabilities:}]
    Geographical coordinates can be associated with buses and network
    elements, but GIS-based transmission-system modelling is not the primary
    focus of OpenDSS. Advanced spatial processing would normally be performed
    using external tools.

    \item[\textbf{Input / output formats:}]
    The native OpenDSS model is based primarily on text-based DSS files and
    commands. Simulation results can be exported for subsequent processing,
    and Python interfaces can be used to retrieve results directly from the
    simulation engine.

    \item[\textbf{Main advantages for EAGLE:}] \hfill
    \begin{itemize}
        \item Mature framework for detailed electrical simulation.
        \item Strong capabilities for time-series analysis.
        \item Detailed modelling of distributed generation, storage and
        voltage-control devices.
        \item Particularly suitable if EAGLE is later extended towards
        distribution-level hosting-capacity studies.
        \item Python automation is possible through the DSS-Extensions
        ecosystem.
    \end{itemize}

    \item[\textbf{Main limitations for EAGLE:}] \hfill
    \begin{itemize}
        \item Primarily designed for distribution and subtransmission
        systems rather than large European transmission-system planning.
        \item Large-scale transmission optimisation is not its main purpose.
        \item Less directly compatible with the PyPSA-Eur data ecosystem.
        \item Integration into a fully Python-native EAGLE workflow requires
        an additional interface such as DSS-Python or OpenDSSDirect.py.
        \item Its detailed unbalanced distribution capabilities may exceed
        the requirements of the initial EAGLE transmission-network model.
    \end{itemize}

    \item[\textbf{Potential application in EAGLE:}]
    OpenDSS would be particularly valuable if EAGLE is extended from the
    high-voltage transmission system towards detailed distribution-network
    hosting-capacity analysis. For the initial national transmission-level
    model, its main strengths are less closely aligned with the current
    EAGLE requirements than those of transmission-oriented Python
    frameworks.

\end{description}


% ============================================================
% MATPOWER
% ============================================================

\subsection{MATPOWER}

\begin{description}

    \item[\textbf{Developer / organisation:}]
    MATPOWER project, originally developed at Cornell University and
    maintained by the MATPOWER development community.

    \item[\textbf{Access:}] \hfill
    \begin{itemize}
        \item Official website:
        \url{https://matpower.org/}

        \item Documentation:
        \url{https://matpower.org/doc/}

        \item Source repository:
        \url{https://github.com/MATPOWER/matpower}
    \end{itemize}

    \item[\textbf{Description:}]
    MATPOWER is a free and open-source package for steady-state electric
    power-system simulation and optimisation in MATLAB and GNU Octave.
    It is widely used in research and education and provides a compact
    bus-branch representation for power-flow and optimal-power-flow studies.

    \item[\textbf{Main capabilities:}] \hfill
    \begin{itemize}
        \item AC power flow.
        \item DC power flow.
        \item Continuation power flow.
        \item AC optimal power flow.
        \item DC optimal power flow.
        \item Economic dispatch and market-related analyses.
        \item Extensible mathematical optimisation framework.
        \item Multi-period optimisation through MOST.
        \item Large collection of benchmark network cases.
    \end{itemize}

    \item[\textbf{Network components:}] \hfill
    \begin{itemize}
        \item Buses.
        \item Branches representing lines and transformers.
        \item Generators.
        \item Loads represented through bus demand.
        \item Generator costs.
        \item Additional components through extensions and the newer
        MATPOWER architecture.
    \end{itemize}

    \item[\textbf{Power-system analysis:}] \hfill
    \begin{itemize}
        \item Newton-based AC power flow.
        \item DC power flow.
        \item Continuation power flow.
        \item AC optimal power flow.
        \item DC optimal power flow.
        \item Multi-period scheduling and optimisation through MOST.
    \end{itemize}

    \item[\textbf{Data handling and automation:}]
    MATPOWER uses MATLAB/Octave structures and M-files to represent electrical
    networks. This provides a compact and highly reproducible representation
    and makes automated studies straightforward within MATLAB or GNU Octave.
    Integration with a Python-centred workflow, however, normally requires
    conversion or an additional interface.

    \item[\textbf{GIS / spatial capabilities:}]
    MATPOWER is primarily an electrical and mathematical simulation framework
    and does not provide a GIS-oriented network data model as a core feature.
    Geographical processing would generally be performed externally.

    \item[\textbf{Input / output formats:}]
    The standard MATPOWER case format stores bus, generator, branch and cost
    information in MATLAB structures. The format is widely supported by other
    open-source power-system tools, making MATPOWER cases useful for network
    interchange and benchmark studies.

    \item[\textbf{Main advantages for EAGLE:}] \hfill
    \begin{itemize}
        \item Mature and extensively used power-flow and OPF framework.
        \item Simple and transparent mathematical network representation.
        \item Large collection of standard test systems.
        \item Useful as an independent reference for validating EAGLE
        power-flow or optimisation results.
        \item MATPOWER case format is supported by several other
        power-system tools.
        \item GNU Octave provides a fully open-source execution alternative
        to MATLAB.
    \end{itemize}

    \item[\textbf{Main limitations for EAGLE:}] \hfill
    \begin{itemize}
        \item Not Python-native.
        \item Limited native GIS and spatial-analysis capabilities.
        \item Less direct integration with the current PyPSA-Eur data
        workflow.
        \item The traditional bus-branch data model is less convenient than
        element-based models when handling detailed equipment information.
        \item Long-term energy-system and renewable-planning functionality
        requires additional tools such as MOST or custom implementations.
    \end{itemize}

    \item[\textbf{Potential application in EAGLE:}]
    MATPOWER could provide an independent reference implementation for
    validating power-flow and optimal-power-flow calculations performed by
    the main EAGLE framework. Its widespread case format may also be useful
    for exchanging network models with other tools. However, adopting it as
    the central EAGLE framework would require additional integration with
    the Python and geospatial workflow.

\end{description}


% ============================================================
% VeraGrid / GridCal
% ============================================================

\subsection{VeraGrid (formerly GridCal)}

\begin{description}

    \item[\textbf{Developer / organisation:}]
    VeraGrid open-source project, led by Santiago Peñate Vera. VeraGrid is
    the current continuation of the project previously known as GridCal.

    \item[\textbf{Access:}] \hfill
    \begin{itemize}
        \item Documentation:
        \url{https://veragrid.readthedocs.io/}

        \item Source repository:
        \url{https://github.com/SanPen/VeraGrid}

        \item Python package:
        \url{https://pypi.org/project/VeraGrid/}
    \end{itemize}

    \item[\textbf{Description:}]
    VeraGrid is an open-source, cross-platform power-system planning and
    simulation platform written in Python. It combines a graphical user
    interface with a Python calculation engine and provides a broad range
    of power-system analysis, optimisation and data-conversion capabilities.
    It is the current continuation of GridCal.

    \item[\textbf{Main capabilities:}] \hfill
    \begin{itemize}
        \item AC/DC multi-grid power flow.
        \item Three-phase unbalanced power flow.
        \item Linear optimal power flow.
        \item PTDF and LODF linear analysis.
        \item Net transfer capacity calculations.
        \item Contingency analysis.
        \item Stochastic power flow.
        \item Short-circuit calculations.
        \item Continuation power flow.
        \item Time-series simulations.
        \item Investment analysis.
        \item RMS and EMT simulation capabilities.
        \item Graphical network editing and visualisation.
        \item Network model validation and repair tools.
    \end{itemize}

    \item[\textbf{Network components:}] \hfill
    \begin{itemize}
        \item Buses and substations.
        \item AC transmission lines.
        \item HVDC connections.
        \item Transformers.
        \item Loads.
        \item Generators.
        \item Batteries and storage elements.
        \item Shunts.
        \item Switches and other network devices.
        \item Detailed line and equipment templates.
    \end{itemize}

    \item[\textbf{Power-system analysis:}] \hfill
    \begin{itemize}
        \item Non-linear AC/DC power flow.
        \item Linearised network analysis.
        \item Linear optimal power flow.
        \item Contingency and N-1 analysis.
        \item PTDF and LODF calculation.
        \item Optimal net transfer capacity analysis.
        \item Stochastic power flow.
        \item Short-circuit and continuation power-flow analysis.
        \item Time-series simulation for multiple study types.
    \end{itemize}

    \item[\textbf{Data handling and automation:}]
    VeraGrid is implemented in Python and separates its calculation engine
    from its graphical interface. The VeraGridEngine package can therefore
    be used directly as a Python library for scripting, automated batch
    calculations and integration into external applications. A server
    package is also available for remote access to the calculation engine.

    \item[\textbf{GIS / spatial capabilities:}]
    VeraGrid includes network schematic and substation-line map
    representations and can handle geographical information associated with
    network assets. Its Python architecture also allows integration with
    external geospatial-processing libraries.

    \item[\textbf{Input / output formats:}]
    VeraGrid supports a broad range of power-system formats, including
    MATPOWER, PSS/E RAW/RAWX, CIM, CGMES, PyPSA, JSON and its native
    VeraGrid formats. This makes it particularly useful for exchanging and
    validating network models between different software ecosystems.

    \item[\textbf{Main advantages for EAGLE:}] \hfill
    \begin{itemize}
        \item Native Python implementation.
        \item Graphical interface in addition to programmatic access.
        \item Broad range of transmission-system analysis functions.
        \item Contingency, PTDF, LODF and transfer-capacity analyses are
        directly relevant to transmission-network constraints.
        \item Supports time-series simulations.
        \item Can import PyPSA network models.
        \item Strong interoperability through MATPOWER, CIM/CGMES, PSS/E
        and other network formats.
        \item Could provide an independent environment for visually
        inspecting and validating the EAGLE network.
    \end{itemize}

    \item[\textbf{Main limitations for EAGLE:}] \hfill
    \begin{itemize}
        \item The current project has recently transitioned from the GridCal
        name to VeraGrid, so older publications and documentation may use
        different package and class names.
        \item Its ecosystem and user community are smaller than those of
        some more established frameworks.
        \item Direct use of PyPSA-Eur requires an import/conversion step,
        even though PyPSA models are supported.
        \item Long-term European energy-system capacity-expansion workflows
        are less directly aligned with the PyPSA-Eur ecosystem.
    \end{itemize}

    \item[\textbf{Potential application in EAGLE:}]
    VeraGrid could be used as an alternative or complementary transmission
    analysis framework for EAGLE. Its Python engine, graphical interface,
    contingency-analysis capabilities and support for PyPSA, CGMES and
    MATPOWER formats make it particularly interesting for inspecting and
    independently validating the Spanish network model and for analysing
    transmission constraints and transfer capacities.

\end{description}